% !TEX root = ../graph-theory.tex
% =====================================================================
%  组合优化 · 图的遍历
% =====================================================================
\lecture{图的遍历：通用的扫描算法及其正确性，深度优先搜索（栈），广度优先搜索（队列），邻接矩阵与邻接表下的实现与复杂度，连通分支，二部图的判定与奇圈的找出}
\label{lec:traversal}

遍历是从一个顶点 $s$ 出发，沿着边把能到达的顶点都访问一遍。几乎所有图算法都以遍历为骨架：求连通分支、判断二部图、求强连通分支、拓扑排序、找欧拉环游、求 st-编号，都是在遍历的过程中顺便记录一些信息。所以这一讲把遍历讲透：先讲一个通用的算法并证明它对，然后看“从哪里继续”的两种选法（栈与队列）分别得到什么，最后分析它在两种存储结构下要花多少时间。

\topic{一、通用的扫描算法}

算法维护三样东西（讲义算法 2.1）：
\begin{itemize}
\item $R$：已经到达的顶点；
\item $Q$：已经到达、但可能还有边没检查过的顶点（“待处理”）；
\item $T$：搜索树的边，每个新顶点 $w$ 进入 $R$ 时，记下是从哪个顶点 $v$ 走过来的。
\end{itemize}

\begin{enumerate}
\item 开始：$R=Q=\{s\}$，$T=\varnothing$。
\item 若 $Q=\varnothing$，停止。否则从 $Q$ 中\emph{选}一个顶点 $v$。
\item 若 $v$ 有不在 $R$ 中的邻点 $w$：把 $w$ 放进 $R$ 和 $Q$，把边 $vw$ 放进 $T$。否则把 $v$ 从 $Q$ 中删去。回到第 2 步。
\end{enumerate}

\begin{theorem}\label{thm:traversal:scan}
无论第 2 步怎样选 $v$，算法都会停止；停止时 $R$ 恰是从 $s$ 可到达的顶点集，$(R,T)$ 是一棵树（有向图中是以 $s$ 为根的分支）。
\end{theorem}
\begin{proof}
\emph{停止。}第 3 步每执行一次，或者 $R$ 增加一个顶点（至多 $n-1$ 次），或者 $Q$ 减少一个顶点；每个顶点只进 $Q$ 一次（进 $Q$ 时同时进 $R$，而进 $R$ 的顶点不再进），所以也只出 $Q$ 一次。第 3 步至多执行 $2n$ 次。

\emph{$(R,T)$ 是树。}每加入一个新顶点 $w$，恰好同时加入一条边 $vw$，$v$ 已在 $R$ 中。所以始终有 $\abs T=\abs R-1$，且 $(R,T)$ 连通；由树的等价刻画（讲义第 1 章定理 1.8），它是树。有向图中每个非 $s$ 的顶点恰有一条进入它的树弧，是分支。

\emph{$R$ 是可达集。}$R$ 中的顶点显然可达。设 $x$ 可达但 $x\notin R$。取一条 $(s,x)$-路，它从 $R$ 出发，最终离开 $R$，所以有一条边 $vw$，$v\in R$，$w\notin R$。$v$ 进过 $Q$，而 $v$ 离开 $Q$ 的条件是 $v$ 没有不在 $R$ 中的邻点；但 $w$ 一直不在 $R$ 中，矛盾。
\end{proof}

第 2 步“选哪个 $v$”留给我们决定。两种选法：
\begin{itemize}
\item 总选\emph{最后}进入 $Q$ 的顶点：$Q$ 是栈，得到\term{深度优先搜索}（DFS）。它一有新顶点就从新顶点继续往前走，走不动了才退回来。
\item 总选\emph{最先}进入 $Q$ 的顶点：$Q$ 是队列，得到\term{广度优先搜索}（BFS）。它先把 $s$ 的邻点都找到，再找邻点的邻点，一层一层地往外推。
\end{itemize}

两种算法都用一个双亲数组 \texttt{tree} 同时记录 $R$ 与 $T$：\texttt{tree[v]}$=-2$ 表示 $v\notin R$，$-1$ 表示 $v=s$ 是根，否则 \texttt{tree[v]} 是 $v$ 在搜索树中的双亲（即边 $(\texttt{tree[v]},v)\in T$）。

\topic{二、深度优先搜索}

“$v$ 有没有不在 $R$ 中的邻点”不能每次都从 \texttt{adj[v]} 的开头重新找，否则同一条边会被反复检查。所以给每个顶点一个下标 \texttt{cur[v]}，表示 \texttt{adj[v]} 中前 \texttt{cur[v]} 个邻点已经看过了，下次从第 \texttt{cur[v]} 个接着看。讲义中邻接表版的 DFS 把 \texttt{visited[i]} 设为边表的头指针、每看一个邻点就令 \texttt{visited[i]=visited[i].next}，是同一个意思。

\begin{lstlisting}[language=Python]
def dfs(adj, s):
    """从 s 出发的深度优先搜索，返回双亲数组 tree。"""
    n = len(adj)
    tree = [-2] * n          # -2：还没到达
    cur = [0] * n            # cur[v]：adj[v] 中下一个要看的位置
    tree[s] = -1             # s 是根
    stack = [s]
    while stack:
        u = stack[-1]                    # 看栈顶，不取走
        if cur[u] < len(adj[u]):         # u 还有没看过的邻点
            v = adj[u][cur[u]]
            cur[u] += 1
            if tree[v] == -2:            # v 是新顶点
                tree[v] = u
                stack.append(v)          # 从 v 继续往前走
        else:                            # u 的邻点都看过了
            stack.pop()                  # 退回到 u 的上一个顶点
    return tree
\end{lstlisting}

\begin{example}\label{ex:traversal:dfs}
在上一讲的图 $G$ 上从 $s=0$ 做 DFS，\texttt{adj=[[1,2,3],[0,2,4],[0,1,3,4],[0,2,4],[1,2,3]]}。下表每一行是 \texttt{while} 循环的一轮：

\begin{center}
\small
\begin{tabular}{rllll}
\toprule
轮 & 栈顶 $u$ & 这一轮看的邻点 & 动作 & 之后的栈 \\
\midrule
0 & & & 初始化，\texttt{tree=[-1,-2,-2,-2,-2]} & \texttt{[0]} \\
1 & 0 & \texttt{adj[0][0]=1} & 新顶点，\texttt{tree[1]=0}，进栈 & \texttt{[0,1]} \\
2 & 1 & \texttt{adj[1][0]=0} & 已到达，跳过 & \texttt{[0,1]} \\
3 & 1 & \texttt{adj[1][1]=2} & 新顶点，\texttt{tree[2]=1}，进栈 & \texttt{[0,1,2]} \\
4 & 2 & \texttt{adj[2][0]=0} & 已到达 & \texttt{[0,1,2]} \\
5 & 2 & \texttt{adj[2][1]=1} & 已到达 & \texttt{[0,1,2]} \\
6 & 2 & \texttt{adj[2][2]=3} & 新顶点，\texttt{tree[3]=2}，进栈 & \texttt{[0,1,2,3]} \\
7 & 3 & \texttt{adj[3][0]=0} & 已到达 & \texttt{[0,1,2,3]} \\
8 & 3 & \texttt{adj[3][1]=2} & 已到达 & \texttt{[0,1,2,3]} \\
9 & 3 & \texttt{adj[3][2]=4} & 新顶点，\texttt{tree[4]=3}，进栈 & \texttt{[0,1,2,3,4]} \\
10 & 4 & \texttt{adj[4][0]=1} & 已到达 & \texttt{[0,1,2,3,4]} \\
11 & 4 & \texttt{adj[4][1]=2} & 已到达 & \texttt{[0,1,2,3,4]} \\
12 & 4 & \texttt{adj[4][2]=3} & 已到达 & \texttt{[0,1,2,3,4]} \\
13 & 4 & （\texttt{cur[4]=3}，看完了） & 出栈 & \texttt{[0,1,2,3]} \\
14 & 3 & （\texttt{cur[3]=3}，看完了） & 出栈 & \texttt{[0,1,2]} \\
15 & 2 & \texttt{adj[2][3]=4} & 已到达 & \texttt{[0,1,2]} \\
16 & 2 & （看完了） & 出栈 & \texttt{[0,1]} \\
17 & 1 & \texttt{adj[1][2]=4} & 已到达 & \texttt{[0,1]} \\
18 & 1 & （看完了） & 出栈 & \texttt{[0]} \\
19 & 0 & \texttt{adj[0][1]=2} & 已到达 & \texttt{[0]} \\
20 & 0 & \texttt{adj[0][2]=3} & 已到达 & \texttt{[0]} \\
21 & 0 & （看完了） & 出栈 & \texttt{[]} \\
\bottomrule
\end{tabular}
\end{center}

结果 \texttt{tree=[-1,0,1,2,3]}：DFS 树是路 $0\,1\,2\,3\,4$。顶点出栈的顺序是 $4,3,2,1,0$，下一讲求强连通分支要用到这个顺序。每个邻接表中的位置恰好被看了一次（$16=2m$ 轮“看邻点”），每个顶点恰好出栈一次（$5=n$ 轮“出栈”），共 $2m+n=21$ 轮。
\end{example}

\begin{center}
\begin{tikzpicture}[scale=0.95]
  \begin{scope}
    \node[gv] (0) at (0,1.5) {0}; \node[gv] (1) at (-1.5,0) {1}; \node[gv] (2) at (0,0) {2};
    \node[gv] (3) at (1.5,0) {3}; \node[gv] (4) at (0,-1.5) {4};
    \foreach \a/\b in {0/2,0/3,1/4,2/4} \draw[ge,black!35] (\a) -- (\b);
    \foreach \a/\b in {0/1,1/2,2/3,3/4} \draw[geh] (\a) -- (\b);
    \node at (0,-2.3) {DFS 树（粗线）};
  \end{scope}
  \begin{scope}[xshift=5cm]
    \node[gv] (0) at (0,1.5) {0}; \node[gv] (1) at (-1.5,0) {1}; \node[gv] (2) at (0,0) {2};
    \node[gv] (3) at (1.5,0) {3}; \node[gv] (4) at (0,-1.5) {4};
    \foreach \a/\b in {1/2,2/3,2/4,3/4} \draw[ge,black!35] (\a) -- (\b);
    \foreach \a/\b in {0/1,0/2,0/3,1/4} \draw[geh] (\a) -- (\b);
    \node at (0,-2.3) {BFS 树（粗线）};
  \end{scope}
\end{tikzpicture}
\end{center}

\topic{三、广度优先搜索}

BFS 中，被选中的顶点是队首。因为队首要一直等到它的所有邻点都处理完才离开，我们可以把“选队首、看一个邻点”的许多轮合并成一轮：取出队首 $u$，一次看完它的所有邻点。这就是讲义算法 2.5 的写法。

\begin{lstlisting}[language=Python]
from collections import deque

def bfs(adj, s):
    """从 s 出发的广度优先搜索，返回双亲数组 tree 与距离 dist。"""
    n = len(adj)
    tree = [-2] * n
    dist = [None] * n
    tree[s], dist[s] = -1, 0
    queue = deque([s])
    while queue:
        u = queue.popleft()               # 取出最早进队的顶点
        for v in adj[u]:                  # 一次看完 u 的所有邻点
            if tree[v] == -2:
                tree[v] = u
                dist[v] = dist[u] + 1
                queue.append(v)
    return tree, dist
\end{lstlisting}

\begin{example}\label{ex:traversal:bfs}
同一个图，从 $s=0$ 做 BFS：

\begin{center}
\small
\begin{tabular}{rlll}
\toprule
轮 & 出队 $u$ & 看 \texttt{adj[u]} 的结果 & 之后的队列 \\
\midrule
0 & & 初始化，\texttt{tree=[-1,-2,-2,-2,-2]} & \texttt{[0]} \\
1 & 0 & $1$ 新（\texttt{tree[1]=0}），$2$ 新（\texttt{tree[2]=0}），$3$ 新（\texttt{tree[3]=0}） & \texttt{[1,2,3]} \\
2 & 1 & $0$ 已到达，$2$ 已到达，$4$ 新（\texttt{tree[4]=1}） & \texttt{[2,3,4]} \\
3 & 2 & $0,1,3,4$ 都已到达 & \texttt{[3,4]} \\
4 & 3 & $0,2,4$ 都已到达 & \texttt{[4]} \\
5 & 4 & $1,2,3$ 都已到达 & \texttt{[]} \\
\bottomrule
\end{tabular}
\end{center}

结果 \texttt{tree=[-1,0,0,0,1]}，\texttt{dist=[0,1,1,1,2]}。与 DFS 树（一条长为 $4$ 的路）不同，BFS 树中每个顶点的深度就是它到 $s$ 的距离（第一部分第~\ref{lec:distance-layers} 讲命题~\ref{prop:distance-layers:bfs}）。
\end{example}

BFS 中顶点按到 $s$ 的距离从小到大出队：队列里任何时候只有距离为 $t$ 与 $t+1$ 的顶点，并且前者都在后者前面。这就是 BFS 能求最短路（每条边长度为 $1$ 时）的原因。

\topic{四、复杂度}

\begin{theorem}\label{thm:traversal:time}
用邻接表存图时，DFS 与 BFS 都在 $O(n+m)$ 时间内完成；用邻接矩阵存图时都是 $O(n^2)$。
\end{theorem}
\begin{proof}
邻接表：每个顶点进栈（队）、出栈（队）各一次，$O(n)$。DFS 中每个位置 \texttt{adj[u][cur[u]]} 被看一次，\texttt{cur[u]} 只增不减；BFS 中每个 \texttt{adj[u]} 在 $u$ 出队时被完整扫描一次。所以看邻点的总次数是 $\sum_ud(u)=2m$（有向图为 $m$）。合计 $O(n+m)$。

邻接矩阵：列出 $u$ 的邻点要扫描整行，$n$ 格。每个顶点的行被扫描一次（DFS 中分几次扫完，但每格只看一次），共 $n^2$ 格。
\end{proof}

矩阵版 DFS（讲义算法 2.2）与表版的结构完全相同，只是把“邻点表中的位置”换成“矩阵这一行的列号”：\texttt{visited[i]}$=k$ 表示第 $i$ 行的前 $k$ 列已经看过。

\begin{lstlisting}[language=Python]
def dfs_matrix(A, s):
    n = len(A)
    tree = [-2] * n
    col = [0] * n                    # 讲义的 visited：第 i 行下次从第 col[i] 列看起
    tree[s] = -1
    stack = [s]
    while stack:
        i = stack[-1]
        while col[i] < n and not (A[i][col[i]] == 1 and tree[col[i]] == -2):
            col[i] += 1              # 跳过不是边、或已到达的列
        if col[i] < n:               # 找到新顶点 j = col[i]
            j = col[i]
            col[i] += 1
            tree[j] = i
            stack.append(j)
        else:
            stack.pop()
    return tree
\end{lstlisting}

在 $G$ 的邻接矩阵上：栈顶 $0$，从第 $0$ 列找起，第 $1$ 列是新顶点，\texttt{col[0]=2}，栈 \texttt{[0,1]}；栈顶 $1$，第 $0$ 列已到达，第 $1$ 列不是边，第 $2$ 列是新顶点，\texttt{col[1]=3}，栈 \texttt{[0,1,2]}；栈顶 $2$，第 $0,1$ 列已到达，第 $2$ 列不是边，第 $3$ 列新，\texttt{col[2]=4}；栈顶 $3$，第 $4$ 列新，\texttt{col[3]=5}；栈顶 $4$，第 $0$ 到 $4$ 列都不行，出栈；之后 $3,2,1,0$ 各自从记下的列号接着找，都找不到，依次出栈。\texttt{tree=[-1,0,1,2,3]}，与表版结果相同（因为两者都按顶点编号从小到大看邻点）。

\topic{五、应用：连通分支}

对每个还没到达的顶点各做一次遍历，每次到达的顶点构成一个连通分支。

\begin{lstlisting}[language=Python]
def components(adj):
    """comp[v] = v 所在连通分支的编号 0, 1, 2, ..."""
    n = len(adj)
    comp = [-1] * n
    c = 0
    for s in range(n):
        if comp[s] == -1:                 # s 还没被任何一次遍历到达
            comp[s] = c
            queue = deque([s])
            while queue:
                u = queue.popleft()
                for v in adj[u]:
                    if comp[v] == -1:
                        comp[v] = c
                        queue.append(v)
            c += 1
    return comp
\end{lstlisting}

外层循环 $n$ 次，所有遍历合起来每个顶点出队一次、每个邻接表扫一次，仍是 $O(n+m)$。

\topic{六、应用：二部图的判定（讲义定理 2.12）}

由第一部分第~\ref{lec:distance-layers} 讲定理~\ref{thm:distance-layers:bipartite} 的证明：在每个连通分支中从某个顶点做 BFS，按 $\texttt{dist}$ 的奇偶分成两部分；若没有边连接同一部分的两点，就是二部划分；若有这样的边 $xy$，由 BFS 的性质 $\texttt{dist}[x]=\texttt{dist}[y]$（相邻顶点的距离至多差 $1$，奇偶相同只能相等），沿双亲从 $x$、$y$ 同时往上走，第一次走到同一个顶点 $z$ 时，两段路加上 $xy$ 就是一个奇圈。

\begin{lstlisting}[language=Python]
def bipartite_or_odd_cycle(adj):
    """返回 ('bipartite', side) 或 ('odd cycle', 圈上顶点的列表)。"""
    n = len(adj)
    tree, dist = [-2] * n, [None] * n
    for s in range(n):                       # 每个分支做一次 BFS
        if dist[s] is None:
            tree[s], dist[s] = -1, 0
            queue = deque([s])
            while queue:
                u = queue.popleft()
                for v in adj[u]:
                    if dist[v] is None:
                        tree[v], dist[v] = u, dist[u] + 1
                        queue.append(v)
    for x in range(n):
        for y in adj[x]:
            if dist[x] == dist[y]:           # 同一层的边
                left, right = [x], [y]
                while left[-1] != right[-1]: # 同时往上走，直到相遇
                    left.append(tree[left[-1]])
                    right.append(tree[right[-1]])
                # left = x ... z，right = y ... z；圈 = x ... z ... y
                return 'odd cycle', left + right[-2::-1]
    return 'bipartite', [d % 2 for d in dist]
\end{lstlisting}

两段路同时往上走时总在同一层，所以会在同一时刻到达公共的祖先 $z$。圈长 $2(\texttt{dist}[x]-\texttt{dist}[z])+1$ 是奇数，两段路除 $z$ 外不交（在 $z$ 之前它们一直处在同一层的不同顶点上），所以确实是圈。

\begin{example}
\texttt{H\_edges=[(0,1),(0,2),(1,3),(2,4),(3,4),(3,5),(4,5),(5,6)]}，从 $0$ 做 BFS：\texttt{dist=[0,1,1,2,2,3,4]}，\texttt{tree=[-1,0,0,1,2,3,5]}。扫描边时，$x=3$ 的邻点依次是 $1,4,5$，$\texttt{dist}[3]=\texttt{dist}[4]=2$。往上走：\texttt{left}$=[3]$，\texttt{right}$=[4]$；$3\neq4$，各走一步：\texttt{left}$=[3,1]$，\texttt{right}$=[4,2]$；$1\neq2$，再走一步：\texttt{left}$=[3,1,0]$，\texttt{right}$=[4,2,0]$；相遇于 $z=0$。\texttt{right[-2::-1]} 是 \texttt{right} 去掉最后一个元素再倒过来，即 $[2,4]$。返回奇圈 $[3,1,0,2,4]$，即 $3\,1\,0\,2\,4\,3$，长为 $5$。

去掉边 $(3,4)$ 后，没有同层的边（$3$ 与 $5$、$4$ 与 $5$ 分别在第 $2$、$3$ 层），返回二部划分 \texttt{[0,1,1,0,0,1,0]}：$\{0,3,4,6\}$ 与 $\{1,2,5\}$。
\end{example}

\begin{center}
\begin{tikzpicture}[scale=0.95]
  \node[gv] (0) at (0,0) {0}; \node[gv] (1) at (1.4,0.9) {1}; \node[gv] (2) at (1.4,-0.9) {2};
  \node[gv] (3) at (2.8,0.9) {3}; \node[gv] (4) at (2.8,-0.9) {4}; \node[gv] (5) at (4.2,0) {5}; \node[gv] (6) at (5.6,0) {6};
  \foreach \a/\b in {3/5,4/5,5/6} \draw[ge] (\a) -- (\b);
  \foreach \a/\b in {0/1,0/2,1/3,2/4,3/4} \draw[geh] (\a) -- (\b);
  \foreach \x/\d in {0/0,1.4/1,2.8/2,4.2/3,5.6/4} \node[idx] at (\x,-1.7) {dist=\d};
\end{tikzpicture}
\end{center}

\topic{小结}

\begin{itemize}
\item 通用扫描算法对任何选法都正确（定理~\ref{thm:traversal:scan}）；选法只影响树的形状。
\item 栈给出 DFS：一路走到底再回退；队列给出 BFS：按距离一层一层地扩展，BFS 树给出最短路。
\item 邻接表下两者都是 $O(n+m)$；关键是每个邻点只看一次（DFS 用 \texttt{cur} 记住看到哪里）。
\end{itemize}

\homework{遍历的练习}[遍历]

\exercisepart{练习}

\begin{exercises}{A}{遍历}
\item\label{ex:traversal:A1} 在 \texttt{H\_edges} 的图上从顶点 $6$ 做 DFS 与 BFS（邻接表按 \texttt{adjacency\_list} 的方式建立），写出 \texttt{tree} 与顶点出栈的顺序。
\begin{solution}
邻接表 \texttt{adj=[[1,2],[0,3],[0,4],[1,4,5],[2,3,5],[3,4,6],[5]]}。

DFS：栈 \texttt{[6]}；$6$ 看 $5$，新，栈 \texttt{[6,5]}；$5$ 看 $3$，新，栈 \texttt{[6,5,3]}；$3$ 看 $1$，新，栈 \texttt{[6,5,3,1]}；$1$ 看 $0$，新，栈 \texttt{[6,5,3,1,0]}；$0$ 看 $1$（已到达）、$2$，新，栈 \texttt{[6,5,3,1,0,2]}；$2$ 看 $0$（已到达）、$4$，新，栈 \texttt{[\dots,2,4]}；$4$ 看 $2,3,5$ 都已到达，出栈 $4$；$2$ 看完，出栈；$0$ 看完，出栈；$1$ 看 $3$（已到达），出栈；$3$ 看 $4,5$ 都已到达，出栈；$5$ 看 $4,6$ 都已到达，出栈；$6$ 看完，出栈。\texttt{tree=[1,3,0,5,2,6,-1]}，出栈顺序 $4,2,0,1,3,5,6$。

BFS：出队 $6$，进 $5$；出队 $5$，进 $3,4$；出队 $3$，进 $1$；出队 $4$，进 $2$；出队 $1$，进 $0$；出队 $2$、$0$，无新顶点。\texttt{tree=[1,3,4,5,5,6,-1]}，\texttt{dist=[4,3,3,2,2,1,0]}。
\end{solution}

\item\label{ex:traversal:A2} 证明 BFS 中的断言：任何时刻队列里的顶点，\texttt{dist} 值从队首到队尾不减，且队首与队尾的 \texttt{dist} 至多差 $1$。
\begin{solution}
对操作次数归纳。初始时队列只有 $s$。设某时刻队列为 $v_1,\dots,v_r$，$\texttt{dist}[v_1]\leqslant\dots\leqslant\texttt{dist}[v_r]\leqslant\texttt{dist}[v_1]+1$。出队 $v_1$ 后，剩下的仍满足不减，且 $\texttt{dist}[v_r]\leqslant\texttt{dist}[v_1]+1\leqslant\texttt{dist}[v_2]+1$。接着进队的顶点 $v$ 的 $\texttt{dist}$ 都是 $\texttt{dist}[v_1]+1$，它不小于队中任何元素（$\leqslant\texttt{dist}[v_1]+1$），并且不超过新队首 $v_2$ 的 $\texttt{dist}+1$（$\texttt{dist}[v_2]\geqslant\texttt{dist}[v_1]$）。所以性质保持。
\end{solution}

\item\label{ex:traversal:A3} 写出递归版的 DFS，并说明它与非递归版访问顶点的顺序相同。为什么在顶点很多时 Python 的递归版会出错？
\begin{solution}
\begin{lstlisting}[language=Python]
def dfs_recursive(adj, s, tree=None, u=None):
    if tree is None:
        tree = [-2] * len(adj)
        tree[s], u = -1, s
    for v in adj[u]:
        if tree[v] == -2:
            tree[v] = u
            dfs_recursive(adj, s, tree, v)
    return tree
\end{lstlisting}
非递归版中“栈顶 $u$ 看下一个邻点、遇到新顶点就进栈”，正是递归版中“$u$ 的 \texttt{for} 循环看下一个邻点、遇到新顶点就递归调用”；递归返回相当于出栈，返回后 $u$ 的 \texttt{for} 循环从上次停下的位置继续，相当于 \texttt{cur[u]}。Python 的函数调用也用一个栈（调用栈）保存“每一层停在哪里”，默认最多约 $1000$ 层；一条路形状的图上 DFS 递归深度为 $n$，$n$ 超过 $1000$ 就会报错 \texttt{RecursionError}。所以本书用显式的栈。
\end{solution}

\item\label{ex:traversal:A4} 修改 BFS，在 $O(n+m)$ 时间内求出从 $s$ 到 $t$ 的一条最短路（顶点序列）；若不可达返回 \texttt{None}。
\begin{solution}
做一次 \texttt{bfs}，再从 $t$ 沿 \texttt{tree} 往上走到 $s$：
\begin{lstlisting}[language=Python]
def shortest_path(adj, s, t):
    tree, dist = bfs(adj, s)
    if tree[t] == -2:
        return None
    path = [t]
    while path[-1] != s:
        path.append(tree[path[-1]])
    return path[::-1]                  # 倒过来，从 s 到 t
\end{lstlisting}
往上走的步数是 $\texttt{dist}[t]\leqslant n$。由 BFS 树的性质，得到的路长为 $\texttt{dist}[t]=d(s,t)$。
\end{solution}
\end{exercises}
